% We introduce ontology-to-tools compilation, which turns an ontology (T-Box) into tool APIs and validators so an LLM can build an ontology-aligned knowledge graph directly from text. A preparation agent translates ontologies into ontology-aware tools with runtime validation, exposed via the Model Context Protocol (MCP). An LLM-driven agent then invokes these tools to construct, repair, and validate KG instances and relations from documents using feedback from constraint checks. In a case study on metal–organic polyhedra (MOP) syntheses in The World Avatar, the framework integrates external domain resources (PubChem, CCDC, and RDKit) without rewriting the core workflow. On 30 papers (6,705 annotations), we achieve a micro-averaged F1 of 0.826, demonstrating ontology-consistent curation with reduced manual prompting and query-template engineering.



We introduce \emph{ontology-to-tools compilation} as a proof-of-principle mechanism for coupling large language models (LLMs) with formal domain knowledge. Ontological specifications are compiled into executable tool interfaces, exposed through the Model Context Protocol (MCP), that LLM-based agents use to create and modify knowledge graph instances while enforcing selected semantic constraints during generation rather than only through post-hoc validation. Compilation is treated as a revision-based qualification process: generated tools and prompts undergo static and contract checks, mock A-Box execution, ontology-level reasoning, and bounded transactional patching before runtime use. We evaluate the same mechanism in two contrasting applications: grounded knowledge-graph construction from metal--organic polyhedra synthesis literature and structured interpretation of German thoracic-surgery operative reports. On 30 synthesis papers, graph-recoverable extraction achieves a micro-averaged precision of 0.844, recall of 0.808, and F1 of 0.826 across four output categories. On 30 operative reports, deterministic 79-field graph projections yield 2{,}335 correct case--field decisions out of 2{,}370; counting exact matches, including correct blanks, as TP and each mismatch as one FP and one FN gives precision, recall, and F1 of 0.985. Although the evaluation units are domain-specific and not directly interchangeable, the results show that one compilation and validation pattern can operationalise different ontology structures, document genres, languages, and target semantics, reducing the need for manually engineered domain-specific extraction interfaces.

